% Part: second-order-logic
% Chapter: syntax-and-semantics
% Section: size-of-domain

\documentclass[../../../include/open-logic-section]{subfiles}

\begin{document}

\olfileid{sol}{syn}{siz}

\olsection{لامتناہی تے \usetoken{S}{enumerable}
  \usetoken{P}{domain} دا بیان}

سیٹ~$M$ (ڈیڈیکائنڈ معنیاں وچ) لامتناہی اُدوں تے صرف اُدوں اے جدوں
اک !!a{injective} (اک بہ اک) فنکشن~$f\colon M
\to M$ موجود ہووے جیہڑا !!{surjective} (ہر ہدف قدر لین والا) نہ ہووے؛ یعنی $\ran{f} \neq M$
ہووے۔ پہلے درجے دی منطق وچ اسیں اک اک-رکنی !!{function}~$f$ نوں
ویکھ سکدے آں تے آکھ سکدے آں کہ فنکشن~$\Assign{f}{M}$، جیہڑا
!!a{structure}~$\Struct{M}$ وچ اوس نوں دتا گیا اے، !!{injective}
اے تے $\ran{f} \neq
\Domain{M}$ اے:
\[
\lforall[x][\lforall[y][(\eq[f(x)][f(y)] \to \eq[x][y])]] \land
\lexists[y][\lforall[x][\eq/[y][f(x)]]].
\]
جے $\Struct{M}$ وچ ایہ !!{sentence} پورا اتردا اے، تاں
$\Assign{f}{M}:
\Domain{M} \to \Domain{M}$، !!{injective} اے؛ ایس لئی
$\Domain{M}$ لازماً لامتناہی اے۔ اُلٹے پاسے، جے $\Domain{M}$
لامتناہی اے تے ایس طرح دا فنکشن موجود اے، تاں اسیں
$\Assign{f}{M}$ نوں اوہ فنکشن دے سکدے آں۔ فیر $\Struct{M}$
وچ ایہ !!{sentence} پورا اترے گا۔ پر ایس طریقے لئی ضروری اے کہ ساڈی
زبان وچ اوس مقصد لئی ورتی غیر منطقی علامت~$f$ موجود ہووے۔
دوجے درجے دی منطق وچ اسیں سِدھا آکھ سکدے آں کہ ایہو جیہا فنکشن
\emph{موجود اے}۔ ہُن~$f$ علامت دی لوڑ نہیں رہندی، تے سانوں خالص
دوجے درجے دی منطق دا !!{sentence} ملدا اے:
\[
\fn{Inf} \ident \lexists[u][(\lforall[x][\lforall[y][(\eq[u(x)][u(y)]
      \lif \eq[x][y])]] \land \lexists[y][\lforall[x][\eq/[y][u(x)]]])].
\]
$\Sat{M}{\fn{Inf}}$ اُدوں تے صرف اُدوں جدوں $\Domain{M}$
لامتناہی اے۔ ہُن اسیں $\fn{Fin} \ident \lnot \fn{Inf}$ تعریف
کر سکدے آں؛ $\Sat{M}{\fn{Fin}}$ اُدوں تے صرف اُدوں جدوں
$\Domain{M}$ متناہی اے۔ خالص پہلے درجے دی منطق دا کوئی اک
!!{sentence}، !!{domain} دے لامتناہی ہون نوں بیان نہیں کر سکدا،
بھاویں ایہو جیہے جملیاں دا اک لامتناہی سیٹ ایہہ گل بیان کر سکدا اے۔
خالص پہلے درجے دی منطق دے !!{sentence}s دا کوئی وی سیٹ ایہو
جیہا نہیں جیہڑا کسے !!a{structure} وچ اُدوں تے صرف اُدوں پورا اترے
جدوں اوہدا دائرہ متناہی ہووے۔

\begin{prop}
$\Sat{M}{\fn{Inf}}$ اُدوں تے صرف اُدوں جدوں $\Domain{M}$ لامتناہی اے۔
\end{prop}

\begin{proof}
$\Sat{M}{\fn{Inf}}$ اُدوں تے صرف اُدوں جدوں
  $\Sat{M}{\lforall[x][\lforall[y][(\eq[u(x)][u(y)] \lif \eq[x][y])]]
    \land \lexists[y][\lforall[x][\eq/[y][u(x)]]]}[s]$ کسے~$s$
  لئی درست ہووے۔ جے ایہہ درست اے، تاں $s(u)$ اک !!a{injective}
  فنکشن اے، تے کوئی $y \in \Domain{M}$،~$s(u)$ دی حاصل قدراں
  وچ نہیں آؤندا۔ اُلٹ پاسے، جے اک !!a{injective}
  $f\colon \Domain{M} \to \Domain{M}$ ایہو جیہا اے جیہدے لئی
  $\ran{f}
  \neq \Domain{M}$، تاں $s(u) = f$ رکھ کے ایہو جہی متغیر
  قدر بندی مل جاندی اے۔
\end{proof}

سیٹ $M$، !!{enumerable}
اُدوں اے جدوں اوہدے !!{element}s دی اک گنتی وار فہرست
\[
m_0, m_1, m_2, \dots
\]
موجود ہووے (دہراؤ توں بغیر، پر ممکن اے متناہی ہووے)۔ ایہو جہی
فہرست اُدوں تے صرف اُدوں موجود اے جدوں کوئی !!a{element}
$z \in M$ تے اک فنکشن $f\colon M \to M$ ایہو جیہے ہون کہ
$z$، $f(z)$، $f(f(z))$، \dots،~$M$ دے سارے !!{element}s
ہون۔ جے فہرست موجود اے تاں $z = m_0$ تے
$f(m_k) = m_{k+1}$ (یا $f(m_k) = m_k$ جے $m_k$ فہرست دا آخری
!!{element} اے) مطلوبہ !!{element} تے فنکشن دندے نیں۔
اُلٹ پاسے، جے ایہو جیہے $z$ تے $f$ موجود نیں، تاں
$z$، $f(z)$، $f(f(z))$، \dots،~$M$ دی گنتی وار فہرست اے؛
ایس لئی $M$، !!{enumerable} اے۔ دوجے درجے دی منطق وچ اسیں
$z$ تے~$f$ دے وجود نوں بیان کر سکدے آں، تے ایہو جیہا
!!{sentence} بنا سکدے آں جیہڑا کسے !!a{structure} وچ اُدوں
تے صرف اُدوں درست ہووے جدوں اوہ !!{structure}، !!{enumerable} ہووے:
\[
\fn{Count} \ident
\lexists[z][\lexists[u][\lforall[X][((X(z) \land
      \lforall[x][(X(x) \lif X(u(x)))]) \lif \lforall[x][X(x)])]]]
\]

\begin{prop}
$\Sat{M}{\fn{Count}}$ اُدوں تے صرف اُدوں جدوں $\Domain{M}$،
!!{enumerable} اے۔
\end{prop}

\begin{proof}
فرض کرو $\Domain{M}$، !!{enumerable} اے، تے $m_0$، $m_1$،
\dots، اوہدی گنتی وار فہرست اے۔ دہراؤ ہٹا کے اسی یقین
کر سکدے آں کہ کوئی $m_k$ دو واری
نہ آوے۔ $f(m_k) = m_{k+1}$ رکھو۔ ہُن
$s(z) = m_0$ تے $s(u) = f$ رکھو۔ اسیں وکھاؤندے آں کہ
\[
\Sat{M}{\lforall[X][((X(z) \land \lforall[x][(X(x) \lif X(u(x)))])
    \lif \lforall[x][X(x)])]}[s]
\]
فرض کرو $M \subseteq \Domain{M}$ کوئی وی ذیلی سیٹ اے۔ ایہ وی فرض کرو کہ
$\Sat{M}{(X(z) \land \lforall[x][(X(x) \lif
X(u(x)))])}[\Subst{s}{M}{X}]$۔ تاں $\Subst{s}{M}{X}(z) \in M$
تے جدوں وی $x \in M$ ہووے، $(\Subst{s}{M}{X}(u))(x) \in M$
وی ہووے گا۔ دوجے لفظاں وچ، کیوں جے
$\varAssign{\Subst{s}{M}{X}}{s}{X}$ اے، $m_0 \in M$ اے
تے جے $x \in M$ تاں $f(x) \in M$؛ ایس لئی
$m_0 \in M$، $m_1 = f(m_0) \in M$،
$m_2 = f(f(m_0)) \in M$، وغیرہ۔ لہٰذا $M = \Domain{M}$،
تے ایس لئی
$\Sat{M}{\lforall[x][X(x)]}[\Subst{s}{M}{X}]$۔ کیوں جے
$M \subseteq
\Domain{M}$ کوئی وی ذیلی سیٹ سی، سو:
$\Sat{M}{\fn{Count}}$۔

ہُن فرض کرو کہ $\Sat{M}{\fn{Count}}$، یعنی
\[
\Sat{M}{\lforall[X][((X(z) \land \lforall[x][(X(x) \lif X(u(x)))])
    \lif \lforall[x][X(x)])]}[s]
\]
کسے~$s$ لئی۔ $m = s(z)$ تے $f = s(u)$ رکھو، تے
$M = \{m,
f(m), f(f(m)), \dots\}$ نوں ویکھو۔ ایس طرح تعریف
کیتا~$M$ صاف طور تے !!{enumerable} اے۔ فیر مفروضے دے مطابق
\[
\Sat{M}{(X(z) \land \lforall[x][(X(x) \lif X(u(x)))])
    \lif \lforall[x][X(x)]}[\Subst{s}{M}{X}]
\]
پورا اتردا اے۔ نال ای $\Sat{M}{X(z)}[\Subst{s}{M}{X}]$ وی
پورا اتردا اے، کیوں جے $M \ni m =
\Subst{s}{M}{X}(z)$ اے؛ تے
$\Sat{M}{\lforall[x][(X(x) \lif
X(u(x)))]}[\Subst{s}{M}{X}]$ وی پورا اتردا اے، کیوں جے
جدوں وی $x \in M$ ہووے، $f(x) \in
M$ وی ہوندا اے۔ سو کیوں جے شرط تے پورا شرطیہ دوویں پورے اتردے نیں، اوس دا
نتیجہ وی پورا اترنا لازمی اے:
$\Sat{M}{\lforall[x][X(x)]}[\Subst{s}{M}{X}]$۔ پر ایس
دا مطلب $M = \Domain{M}$ اے، تے ایس لئی $\Domain{M}$،
!!{enumerable} اے، کیوں جے تعریف مطابق $M$ اے۔
\end{proof}

\begin{prob}
!!{sentence} $\fn{Inf} \land \fn{Count}$ عین اوہناں دائریاں وچ
درست اے جیہڑے !!{denumerable} نیں۔ $\fn{Count}$ دی تعریف
ایویں بدلو کہ اک وکھرا !!{sentence} بنے جیہڑا دائرے دے
!!{denumerable} ہون نوں سِدھا بیان کرے، تے ثابت کرو کہ اوہ
ایویں ہی کردا اے۔
\end{prob}

\end{document}
